// تفاصيل دورة الرواتب — مكوّن خادم
import { createClient } from '@/lib/supabase-server'
import { redirect, notFound } from 'next/navigation'
import Link from 'next/link'
import RunActions from './RunActions'
import { isStaff, type Role } from '@/lib/roles'

const MONTHS = ['يناير','فبراير','مارس','أبريل','مايو','يونيو',
                'يوليو','أغسطس','سبتمبر','أكتوبر','نوفمبر','ديسمبر']

const STATUS: Record<string, { label: string; bg: string; fg: string }> = {
  draft:     { label: 'مسودة',  bg: '#F1F3F6', fg: '#556' },
  approved:  { label: 'معتمدة', bg: '#E7F0FB', fg: '#1B4F8A' },
  paid:      { label: 'مصروفة', bg: '#E6F6EC', fg: '#1B6B3A' },
  cancelled: { label: 'ملغاة',  bg: '#FBE9E9', fg: '#8A2B2B' },
}

export default async function RunPage({ params }: { params: Promise<{ id: string }> }) {
  const { id } = await params
  const supabase = await createClient()
  const { data: { user } } = await supabase.auth.getUser()
  if (!user) redirect('/login')

  const [{ data: profile }, { data: run }, { data: items }, { data: issues }] = await Promise.all([
    supabase.from('profiles').select('role').eq('id', user.id).single(),
    supabase.from('payroll_runs').select('*').eq('id', id).single(),
    supabase.from('payroll_items').select('*').eq('run_id', id).order('employee_name'),
    supabase.rpc('validate_wps_run', { p_run_id: id }),
  ])

  const role = (profile?.role ?? 'admin') as Role
  if (!isStaff(role)) redirect('/dashboard')
  if (!run) notFound()

  const st = STATUS[run.status] ?? STATUS.draft
  const pasiEmp = (items ?? []).reduce((a, i) => a + Number(i.pasi_employee ?? 0), 0)

  return (
    <div style={{ maxWidth: 1100, margin: '0 auto' }} dir="rtl">
      <Link href="/payroll" style={{ color: '#667', fontSize: 13.5, textDecoration: 'none' }}>
        → عودة لدورات الرواتب
      </Link>

      <div style={{ display: 'flex', justifyContent: 'space-between', alignItems: 'flex-start',
                    flexWrap: 'wrap', gap: 10, marginTop: 10 }}>
        <div>
          <h1 style={{ color: '#0F2744', marginBottom: 6 }}>
            رواتب {MONTHS[run.period_month - 1]} {run.period_year}
          </h1>
          <span style={{ background: st.bg, color: st.fg, padding: '4px 12px',
                         borderRadius: 20, fontSize: 13, fontWeight: 600 }}>
            {st.label}
          </span>
        </div>
        <RunActions
          runId={run.id}
          status={run.status}
          role={role}
          hasIssues={(issues ?? []).length > 0}
        />
      </div>

      {issues && issues.length > 0 && (
        <div style={{ background: '#FBF3D5', border: '1px solid #EAD9A0', borderRadius: 13,
                      padding: 16, margin: '18px 0' }}>
          <b style={{ color: '#7A5C0A' }}>⚠️ بيانات ناقصة تمنع تصدير ملف حماية الأجور</b>
          <ul style={{ color: '#8A6D0F', fontSize: 13.5, margin: '8px 0 0', paddingRight: 18 }}>
            {issues.map((x: any, n: number) => (
              <li key={n}>{x.employee_name} — {x.issue}</li>
            ))}
          </ul>
        </div>
      )}

      <div style={{ marginTop: 18, background: '#0F2744', color: '#fff', borderRadius: 14,
                    padding: 18, display: 'grid',
                    gridTemplateColumns: 'repeat(auto-fit,minmax(150px,1fr))', gap: 14 }}>
        <Sum label="إجمالي الرواتب" v={run.total_gross} />
        <Sum label="اشتراك الموظفين" v={pasiEmp} />
        <Sum label="حصة صاحب العمل" v={run.total_pasi_er} />
        <Sum label="صافي المستحق" v={run.total_net} />
      </div>

      <div style={{ marginTop: 18, background: '#fff', border: '1px solid #E6E9EF',
                    borderRadius: 14, overflowX: 'auto' }}>
        <table style={{ width: '100%', borderCollapse: 'collapse', fontSize: 13.5, minWidth: 760 }}>
          <thead>
            <tr style={{ background: '#F7F9FC', color: '#556' }}>
              <Th>الموظف</Th><Th>الأساسي</Th><Th>البدلات</Th>
              <Th>الخصومات</Th><Th>اشتراك الموظف</Th><Th>حصة صاحب العمل</Th><Th>الصافي</Th>
            </tr>
          </thead>
          <tbody>
            {(items ?? []).map((i) => (
              <tr key={i.id} style={{ borderTop: '1px solid #EEF1F5' }}>
                <Td>{i.employee_name}</Td>
                <Td>{fmt(i.basic_salary)}</Td>
                <Td>{fmt(i.extra_income)}</Td>
                <Td>{fmt(i.deductions)}</Td>
                <Td>{fmt(i.pasi_employee)}</Td>
                <Td>{fmt(i.pasi_employer)}</Td>
                <Td><b>{fmt(i.net_salary)}</b></Td>
              </tr>
            ))}
          </tbody>
        </table>
      </div>
    </div>
  )
}

function Sum({ label, v }: { label: string; v: number | null }) {
  return (
    <div>
      <div style={{ fontSize: 12, opacity: 0.8 }}>{label}</div>
      <div style={{ fontSize: 20, fontWeight: 700, fontFamily: 'Cairo' }}>{fmt(v)}</div>
    </div>
  )
}

function Th({ children }: { children?: React.ReactNode }) {
  return <th style={{ textAlign: 'right', padding: '11px 13px', fontWeight: 600 }}>{children}</th>
}

function Td({ children }: { children?: React.ReactNode }) {
  return <td style={{ padding: '11px 13px', color: '#334' }}>{children}</td>
}

function fmt(v: number | null) {
  return Number(v ?? 0).toLocaleString('en-US', { minimumFractionDigits: 3, maximumFractionDigits: 3 })
}
